\documentclass[sigconf,nonacm,11pt]{acmart}

\usepackage{booktabs}
\usepackage{amsmath}
\usepackage{graphicx}

% colonnes etroites + identifiants longs non coupables
\setlength{\emergencystretch}{3em}

\settopmatter{printacmref=false}
\renewcommand\footnotetextcopyrightpermission[1]{}
\pagestyle{plain}

\begin{document}

\title{Bytes Are Not Milliseconds}
\subtitle{A Measurement Study of Two-Pass Block Selection for Long-Context LLM Inference}

\author{Anonymous}
\affiliation{\institution{}\country{}}

\begin{abstract}
Block-sparse attention makes long-context decoding tractable by reading only a
fraction of the KV cache, but the \emph{selector}---the mechanism that decides
which blocks to read---has itself become the dominant cost. We study
\emph{Adaptive Sequential Precision} (ASP), a two-pass selector that scores all
blocks with a coarse per-block summary, then rescores only the survivors with a
finer one, and we formalise it as a \emph{precision-allocation bandit}: a
best-arm-identification problem in which the measurement budget is spent in
\emph{bits per arm} along a rate--distortion curve rather than in repeated
stochastic pulls.

Our first headline result is negative and, we argue, instructive. On Qwen3-8B, ASP
reduces selector KV traffic by $13\text{--}25\times$ yet is $1.1\text{--}2.5\times$
\emph{slower} end-to-end on two independent backends. We explain the gap in three measured steps. First, the
irregular gather that ASP requires is essentially free on datacenter hardware:
across four memory architectures (LPDDR5, GDDR6, HBM2e, GDDR6X) scattered reads
reach $97\text{--}104\%$ of contiguous bandwidth once each gathered chunk exceeds
$\approx$1\,KiB, and on A100 and RTX 4090 the granularity threshold disappears
entirely.
Second, fusing the two passes into a single kernel with interleaved stratified
selection yields $1.57\times$ over the two-kernel form on an H200 (15/15
configurations), but only when the per-stratum quota is small ($q\!\approx\!8$);
at $q=64$ fusion loses by $3\times$. Third, the residual end-to-end gap is
per-layer kernel-launch overhead: an unfused PyTorch selector costs a flat
$\approx$20--25\,ms per decode step, which alone exceeds the entire dense attention
time ($2\text{--}11$\,ms at $16$k--$262$k tokens on the H200, itself running at
$84\%$ of peak bandwidth; $2.6\text{--}42.7$\,ms at the same lengths on an
RTX 4090, running at $80\text{--}100\%$ of peak). Because the selector's cost is
flat, the crossover moves with the hardware: we bracket it at $N\approx146$k
tokens on the RTX 4090 and extrapolate $N\approx700$k on the H200.

A second, later result revises what we believe about selector \emph{quality}, and
we report it with the same prominence because it changes how the first result
should be read. An audit found that our perplexity harness omitted the causal
mask, so every quality measurement in the first version was contaminated. Under a
corrected harness, the mean-pooled block key that ASP scores with ranks blocks
\emph{worse than chance}, whereas the maximum of the per-key inner products inside
a block matches an exact-mass oracle to within $0.0002$\,nat on SmolLM2-135M and
$0.001$\,nat on Qwen2.5-0.5B. Projecting keys onto
eight principal components per head---$12.5\%$ of key bytes---preserves this to
within $0.011$\,nat ($0.022$ on Qwen2.5-0.5B) when the basis is fit on the first
$256$ keys, needs no training, survives freezing
at prefill, and transfers across texts and domains. The near-tie that we
previously took to cap selector quality is an artefact of the estimator, not of
the problem.

We additionally report two \emph{refuted} predictive models of the gather
granularity threshold---one based on Little's law, one on DRAM page
invariance---and a measurement-driven law relating selection recall to
estimator noise, whose residuals are explained by Jensen's inequality over
per-query noise dispersion ($r=0.96$). We argue that convergent empirical
behaviour across four memory architectures without a reliable predictive model is itself a
result worth documenting.
\end{abstract}

\keywords{LLM inference, sparse attention, KV cache, GPU kernels, memory bandwidth}

\maketitle

%======================================================================
\section{Introduction}

Decoding a token from a long-context language model is bound by memory
bandwidth, not by arithmetic. Every generated token requires reading the entire
key--value (KV) cache, whose size grows linearly with context. At one million
tokens a mid-sized model must stream tens of gigabytes \emph{per token per
sequence}, and no amount of additional compute helps.

Block-sparse attention answers this by reading only the $k$ most relevant blocks
of the cache~\cite{yuan2025nsa,tang2024quest}. This shifts the problem rather
than removing it: to know which $k$ blocks matter, something must look at all
$n$ of them. That \emph{selector} is dense by construction, and it has quietly
become the dominant term. Reconstructing the cost of DeepSeek-V3.2 from its
published hyperparameters, we find that at 1M context its lightning indexer
accounts for $\approx$90\% of decode FLOPs and must stream $7.9$\,GB per decoded
token---while the sparse attention it enables costs $3\%$.\footnote{Derived
independently from the architecture description
in~\cite{deepseek2025v32,deepseek2024v3}, whose MLA latent cache originates
in~\cite{deepseek2024v2}; our reconstruction reproduces the
published parameter counts of DeepSeek-V4 to within $1.7\%$, which we take as
evidence the cost model is faithful.}

\paragraph{Why the existing approach is suboptimal.} Every selector we are aware
of---NSA~\cite{yuan2025nsa}, Quest~\cite{tang2024quest},
COBS~\cite{tian2026cobs}, DeepSeek's CSA~\cite{deepseek2026v4},
AsyncTLS~\cite{hu2026asynctls}---allocates the \emph{same} number of summary
bytes to every block and reads each summary exactly once. This is a uniform,
single-pass design. It is an assumption, not a necessity: most blocks are far
from the decision boundary and could be discarded from a much coarser summary.

\paragraph{Contribution.} We formalise non-uniform, sequential summary reading
as a \emph{precision-allocation bandit}, implement it as a single fused GPU
kernel with stratified selection, and measure it end-to-end on a real model---%
finding that the byte savings are real, large, and \emph{do not become
speedups} unless the selector is fused, and unless the surrounding engine is
itself memory-bound.

%======================================================================
\section{Background and Related Work}

\paragraph{Sparse attention and block selection.} Fixed-pattern methods attend
to predetermined windows and global tokens~\cite{beltagy2020longformer,%
zaheer2020bigbird}. Eviction methods discard low-importance
entries~\cite{xiao2024streamingllm,zhang2023h2o,li2024snapkv}. Query-aware block
selectors summarise contiguous blocks offline and, per query, run fine-grained
attention over the highest-scoring ones; this family includes
Quest~\cite{tang2024quest} and NSA~\cite{yuan2025nsa}. COBS~\cite{tian2026cobs}
formalises selection as ranking blocks by \emph{attention mass} and shows that
first-order (mean-pool) summaries are the limiting factor, proposing a
second-order cumulant correction. Our work is complementary and, on one point,
in tension: we find that at matched byte budget a \emph{support}-based summary
(a coreset of $r$ centroids scored by log-sum-exp) dominates a \emph{moment}-based
one (mean plus rank-$r$ covariance) at equal or lower cost: a coreset of $r=2$
centroids ($2.03$ vectors per block) reaches $86.4\%$ mass recall at $k=8$, where
COBS at $r=8$ ($9.12$ vectors) reaches $81.5\%$. The reason is that the
within-block key covariance is not low-rank---its effective rank is $17.8$ out of $64$ in our measurements, so
a rank-$r$ summary captures only a varying fraction of the quadratic term. The
same measurement bounds what methods built on low-rank key structure, such as
Loki~\cite{singhania2024loki}, can recover \emph{within} a block---though Loki
estimates its rotated basis globally, where the spectrum is more favourable.

\paragraph{Hierarchical selection.} AsyncTLS~\cite{hu2026asynctls} performs
coarse block filtering followed by fine token selection;
HiSparse~\cite{hisparse2026} hierarchises \emph{storage} between host and device,
and lookahead sparse attention~\cite{lsa2026} amortises the index across steps.
Both refine \emph{what} is examined. ASP is orthogonal: it refines \emph{how
precisely} the same objects are examined, and allocates that precision according
to distance from the decision threshold.

\paragraph{Serving systems.} PagedAttention~\cite{kwon2023vllm} and
RadixAttention~\cite{zheng2024sglang} govern KV memory;
FlashAttention-3~\cite{shah2024flashattention3} and
FlashInfer~\cite{ye2025flashinfer} provide the kernels; prefill/decode
disaggregation~\cite{zhong2024distserve,qin2024mooncake,agrawal2024sarathi}
governs placement. ASP sits below all of these, inside the attention kernel.

\paragraph{Speculative decoding.} EAGLE-3~\cite{li2025eagle3} and its
predecessors~\cite{leviathan2023speculative} amortise the KV read across
$\gamma+1$ candidate tokens. We note in passing---and quantify in
Section~\ref{sec:disc}---that sparse attention and speculation interact
negatively, because candidate positions select different blocks and the
verification pass must read their union.

\paragraph{Bandits.} Selecting the top-$k$ of $n$ noisy quantities under a
measurement budget is best-arm identification~\cite{bubeck2009pure,%
carpentier2016tight,karnin2013almost}. We import that framework, with one
structural difference developed in Section~\ref{sec:method}.

%======================================================================
\section{Method}\label{sec:method}

\subsection{Block summarisation as measure quantisation}

A block $b$ is an empirical measure on key space, $\mu_b=\sum_{j\in b}\delta_{k_j}$,
and its attention mass is the Laplace transform of that measure,
$m_b(q)=\int e^{s\,q^\top k}\,d\mu_b(k)$ with $s=1/\sqrt{D}$. Summarising a block
means replacing $\mu_b$ by a few-atom measure whose transform agrees on the query
distribution. Writing $\ell_b=\log m_b$, the Gibbs variational identity gives
\begin{equation}
\ell_b(q)=\max_{p\in\Delta_L}\Big[\textstyle\sum_r p_r\,(s\,q^\top k_r)+H(p)\Big],
\end{equation}
so $\ell_b$ interpolates between $\log L+\text{mean}$ as $q\to 0$ and
$\max_r s\,q^\top k_r$ as $\|q\|\to\infty$. Two approximation families follow:
\emph{moment} approximations expand the cumulant generating function about $q=0$
(mean-pool at order one, COBS at order two), while \emph{support} approximations
approximate the geometry of the key cloud.

\begin{proposition}[Uniform two-sided coreset bound]\label{prop:bound}
Partition a block into clusters $C_1,\dots,C_r$ with centroids $\mu_c$, counts
$n_c$ and radii $\rho_c=\max_{k\in C_c}\|k-\mu_c\|$. Then for \emph{every} query $q$,
\begin{equation}
\Big|\,\ell_b(q)-\log\textstyle\sum_c n_c e^{s\,q^\top\mu_c}\Big|\;\le\;s\,\|q\|\max_c \rho_c .
\end{equation}
\end{proposition}
\begin{proof}
For $k\in C_c$, Cauchy--Schwarz gives $|q^\top(k-\mu_c)|\le\|q\|\rho_c$, hence
$n_c e^{s q^\top\mu_c-s\|q\|\rho_c}\le\sum_{k\in C_c}e^{s q^\top k}\le
 n_c e^{s q^\top\mu_c+s\|q\|\rho_c}$. Summing over $c$ and taking logarithms, the
common factors $e^{\pm s\|q\|\rho}$ factor out.
\end{proof}

Unlike a truncated cumulant expansion, this bound is uniform in the direction of
$q$, is sign-aware (the term $e^{q^\top\mu_c}$ distinguishes $+\mu$ from $-\mu$,
whereas $q^\top\Sigma q$ does not), and is exact on an isolated outlier, which
forms a singleton of zero radius. By Jensen's inequality the coreset estimate is
moreover always a \emph{lower} bound on $\ell_b$, so the interval is one-sided
and half as wide as the proposition suggests. The bound also identifies the right
clustering objective: minimising the \emph{maximum} radius, i.e.\
$k$-center~\cite{gonzalez1985clustering}, rather than $k$-means.

\subsection{The precision-allocation bandit}

Let $n$ blocks be candidates and $k$ be selected. A summary of $b$ bytes yields an
estimator with noise $\sigma(b)$, a measurable rate--distortion curve---the same
curve that governs KV-cache quantization~\cite{turboquant2025,blockgtq2026}, applied
there uniformly to every block. A two-pass
scheme reads $b_1$ bytes for all $n$ blocks and $b_2$ bytes for $m$ survivors:
\begin{equation}
B(b_1,b_2,m)=n\,b_1+m\,b_2 .
\end{equation}
This is best-arm identification with a crucial difference: in a classical bandit,
pulling an arm $t$ times reduces noise as $1/\sqrt t$ by stochastic averaging;
here, reading $b$ bytes reduces it along a \emph{deterministic} rate--distortion
curve. We call this a \emph{precision-allocation bandit}. To our knowledge this
variant is treated neither in the bandit literature nor in the attention literature.

Recall loss decomposes into two first-order-independent sources. A block of true
rank $j\le k$ is wrongly eliminated in pass~1 if its noisy score falls below that
of the rank-$m$ block; both scores being independently perturbed, their difference
has standard deviation $\sigma_1\sqrt2$, giving
\begin{equation}
L_1(b_1,m)=\sum_{j\le k} w_j\,\Phi\!\Big(\!-\Delta_{j,m}\big/(\sigma(b_1)\sqrt2)\Big),
\label{eq:l1}
\end{equation}
with $\Delta_{j,m}=\ell_{(j)}-\ell_{(m)}$ and $w_j$ the mass share of rank $j$.
Pass-2 ranking error is $L_2(b_2)=1-R(\sigma(b_2))$, where $R$ is the empirical
recall-versus-noise curve of Section~\ref{sec:law}. The design problem is
\begin{equation}
\min_{b_1,b_2,m} \; n\,b_1+m\,b_2
\quad\text{s.t.}\quad R(\sigma(b_2))-L_1(b_1,m)\ \ge\ \tau .
\end{equation}
All quantities---$\sigma(b)$, $\Delta_{j,m}$, $w_j$, $R$---are measurable offline
on a calibration sample, so the configuration is computed, not swept.

\subsection{Fusion by interleaved stratification}\label{sec:fusion}

\begin{figure}[t]
\centering
\includegraphics[width=\columnwidth]{figs/strat.pdf}
\caption{Cost of stratified selection relative to the global top-$m$, versus the
number of strata $G$. Interleaved stratification (solid) is nearly free;
contiguous stratification (dashed) is not, because relevance is correlated
between adjacent blocks.}
\label{fig:strat}
\end{figure}


A naive fused kernel must satisfy a global dependency: pass~2 needs the top-$m$
of pass~1 across all blocks of a query. Three designs are possible. One work-group
per query needs no synchronisation but collapses parallelism to the number of
queries (batch $\times$ KV heads), typically two orders of magnitude below GPU
residency. A grid-wide barrier preserves parallelism but costs $5$--$10\,\mu$s on
A100---the same order as the launch it was meant to save.

We take a third route. ASP is a heuristic filter, not an exact ranking, so the
\emph{global} top-$m$ is not required. We partition the $n$ blocks into $G$
interleaved strata; work-group $g$ of a query handles blocks
$\{g,\,g\!+\!G,\,g\!+\!2G,\dots\}$, keeps the top-$q$ of its own stratum with
$q=m/G$, and immediately gathers its own survivors. There is \emph{no
synchronisation at all}, parallelism is complete, and there is a single launch.

Interleaving is not a detail. Measured on real attention maps (Figure~\ref{fig:strat}),
interleaved stratification costs at most $2.9$--$3.2$ relative recall points across
all $(m,G)$, whereas \emph{contiguous} stratification costs up to $22.3$ points:
adjacent blocks are correlated in relevance, so a contiguous stratum can hold a
disproportionate share of the mass and be truncated by its quota. Interleaving
breaks that correlation, and it is also the natural grid-stride access pattern.

%======================================================================
\section{Experimental Setup}

\paragraph{Hardware.} Five GPUs. The gather sweep spans four memory
technologies: an AMD Radeon 860M integrated GPU (\texttt{gfx1152}, RDNA 3.5,
shared LPDDR5, 80\,GB/s measured), an NVIDIA Tesla T4 (GDDR6, 274\,GB/s
measured), an NVIDIA A100 PCIe 40\,GB (HBM2e, 1372\,GB/s measured) and an NVIDIA
RTX 4090 (GDDR6X, 933\,GB/s measured, 72\,MB of L2). An NVIDIA H200 NVL
(143\,GB, 132 SMs, HBM3e) is used for kernel fusion and end-to-end measurement. Kernels are OpenCL on the AMD part and CUDA
(\texttt{nvcc -O3}) elsewhere; a single source compiles under both \texttt{nvcc}
and \texttt{hipcc}, using no warp-width-dependent primitive.

\paragraph{Selection-quality measurements.} Real $Q$/$K$ activations captured from
Qwen2.5-0.5B and SmolLM2-135M on WikiText at 8192 tokens. Following the
convention of block-sparse systems, the local window and attention sink are
served by separate branches and are therefore \emph{excluded} from the candidate
set; without this exclusion the metric is dominated by trivially-retrieved blocks.
The metric is \emph{mass recall}: the true attention mass captured by the $k$
selected blocks, divided by that of the oracle top-$k$.

\paragraph{End-to-end.} Qwen3-8B~\cite{qwen3} (36 layers, 32 query heads, 8 grouped-query
KV heads~\cite{ainslie2023gqa}, head dim 128, 144\,KiB/token in bf16) on the H200 and, independently, on an
RTX 4090 (GDDR6X, 933\,GB/s measured), with ASP registered
through the \texttt{ALL\_ATTENTION\_\allowbreak FUNCTIONS} extension point of
\texttt{transformers}~5.16---the supported hook, not a monkey-patch. The rest of
the model runs unmodified, so the measurement is a \emph{complete decode step}.
On the 4090 the harness uses the same configuration ($L_b=64$, $r_1=2$, $r_2=8$,
$m_{\text{frac}}=8$, $k_{\text{frac}}=64$, $512$ local keys, $64$ sink keys) and a
synthetic KV tensor: the 4090 run measures latency, not quality.

%======================================================================
\section{Results}

\subsection{A law relating recall to estimator noise}\label{sec:law}

\begin{figure}[t]
\centering
\includegraphics[width=\columnwidth]{figs/law.pdf}
\caption{Reference curve $R(\sigma)$, obtained by injecting controlled Gaussian
noise into the \emph{exact} block log-masses and re-ranking. Three real summaries
are placed at their measured $\sigma$: coreset and mean-pool land on the curve,
COBS beats it by 5--8 points.}
\label{fig:law}
\end{figure}


We first establish what any selector can achieve. Injecting controlled Gaussian
noise of standard deviation $\sigma$ into the \emph{exact} block log-masses and
ranking gives a reference curve $R(\sigma)$ (Figure~\ref{fig:law}) that is
independent of any method: $95\%$ mass recall at $k=8$ requires
$\sigma\le0.39$~nats, and $99\%$ requires $\sigma\le0.16$.

Real methods do not all fall on this curve. Coresets and mean-pool fall within $0.9$
points of the prediction; COBS and Gaussian-mixture summaries beat it by $4.8$ to
$8.1$ points. We tested and rejected two explanations---a rank-dependent bias
(injecting the measured bias makes predictions \emph{worse}, because the
rank-conditional mean of the error is a regression-to-the-mean artefact) and
over-counting of between-stratum variance (the effect is negligible). The
explanation that survives is Jensen's inequality over \emph{per-query} noise
dispersion: $R$ is convex, so $\mathbb{E}_i R(\sigma_i)>R(\mathbb{E}_i\sigma_i)$,
and the gap correlates with the coefficient of variation of $\sigma$ across
queries at $r=+0.956$. The law is therefore correct but must be applied per query
and then averaged.

\paragraph{Blocks are nearly tied.} The median log-mass gap between the 8th and
9th best block is $0.070$~nats, against a residual estimator uncertainty of
$1.08$~nats without RoPE and $2.36$ with---a ratio of $15$ to $34$. No compact summary can resolve the ordering at
the cut. This single fact explains four independent negative results we obtained
(quadrature-optimised weights, correctness certificates, learned calibration and
adaptive $k$ all reduce estimation error without improving ranking), and it
explains why mass recall ($85$--$90\%$) so far exceeds set recall ($60$--$75\%$):
mis-selecting a near-tied block is nearly free.

\subsection{The irregular gather is free}

\begin{figure*}[t]
\centering
\includegraphics[width=\textwidth]{figs/gather.pdf}
\caption{Scattered versus contiguous reads at equal volume, on four memory
architectures. \textbf{Left:} bandwidth ratio. \textbf{Right:} fraction of peak
bandwidth reached by the contiguous path; below $\approx$30\% neither path
saturates and the ratio is uninformative. The gather penalty vanishes above
$\approx$1\,KiB and decreases monotonically with GPU parallelism; on A100 and
RTX 4090 the threshold has disappeared.}
\label{fig:gather}
\end{figure*}

ASP's second pass reads a query-dependent subset of summaries, an irregular
gather. The concern is misposed: each gathered element is a whole block summary
of $r\!\cdot\! D\!\cdot\!2$ bytes---1\,KiB for $r=8$, $D=64$, i.e.\ sixteen
64-byte cache lines---so coalescing is perfect \emph{within} a chunk and only
inter-chunk locality is lost.

Figure~\ref{fig:gather} confirms this and locates the threshold. In the saturated
regime (contiguous path above $80\%$ of peak) the gather averages $104.0\%$ of
contiguous bandwidth on the integrated GPU, $98.4\%$ on T4, $97.7\%$ on A100 and
$98.8\%$ on the RTX 4090, and never falls below $94.9\%$ on any of the four;
values above unity are measurement noise on a path that is already
bandwidth-saturated. Below saturation the four separate: the integrated GPU falls
to $67.0\%$ at 64-byte chunks and T4 to $88.0\%$ at 512 bytes, whereas A100 never
drops below $95.5\%$ at \emph{any} granularity and the RTX 4090 never drops below
$96.9\%$ above 512 bytes---the threshold has disappeared on both. On the 4090 the
smallest saturated chunk is 512 bytes, so the design rule
$r_2\!\cdot\!D\!\cdot\!2\ge2048$ bytes is conservative rather than tight.

\paragraph{An efficiency number that measures cache residency.} One statistic
from the gather bench does \emph{not} survive scrutiny, and we report the failure
because it changed how we read our own tables. The bench also reports an
\emph{efficiency}: the measured speedup of the two-pass path over a contiguous
read of the same volume, divided by the byte ratio. On the 4090 it averages
$47.0\%$, and $20$ of $32$ configurations are slower than the flat path. The two
configurations above $200\%$ ($278.4\%$ and $222.3\%$) look like outliers and are
in fact the only honest ones. The flat baseline reads $n\,r_2D\,2$ bytes; on a
4090 that volume exceeds the $72$\,MB L2 only for those two rows. Everywhere else
the flat path is served from L2---its implied bandwidth is $822$--$3647$\,GB/s,
above the $933$\,GB/s DRAM peak measured on the same card---so the ratio measures
cache residency, not DRAM behaviour. The same regime split appears on
every card: efficiency is $28.3\%$ (T4, $4$\,MB L2), $30.2\%$ (A100, $40$\,MB)
and $33.5\%$ (4090, $72$\,MB) whenever the flat volume fits in L2---three GPUs
whose launch costs and bandwidths differ by an order of magnitude and which agree
to five points---and $81.0\%$ (T4) and $75.9\%$ (A100) when it does not. On T4 and
A100 volume and cache regime vary together, so the split cannot be attributed to
the cache alone; on the 4090 it can, because the implied flat bandwidth reaches
$2097$\,GB/s against a measured DRAM peak of $933$\,GB/s. Not one of the $32$ configurations compares two DRAM-bound paths, because
the ASP path never exceeds $50$\,MB. The granularity sweep above is unaffected: it
reads $96$\,MiB per launch, above every L2 in the study. The design rule is
$r_2\!\cdot\! D\!\cdot\!2\ge2048$ bytes. The coarse pass violates it by
construction ($r_1=2$, $D=64$: 256 bytes), which is harmless precisely because
that pass reads every block and is therefore a contiguous scan, not a gather.

\paragraph{Two refuted models.} We pre-registered a predictive model of the
threshold based on Little's law~\cite{little2008queueing}: a gather kernel
saturates when $N_{\text{res}}\cdot\min(C,U)\ge B\lambda$, giving
$C^\star=\min(U,\alpha B\lambda/N_{\text{res}})$ with $\alpha$ calibrated on the
integrated GPU. It predicted $328$~bytes for T4---\emph{lower} than the 1024
observed there---and $1009$ for A100. Measurement refuted the calibration twice:
T4's threshold is $1024$, and A100 has none. We then adopted a DRAM-page
invariance hypothesis (the threshold is $\approx$1\,KiB on any memory system);
A100 refutes that too. What survives is only the \emph{direction} of the first
model: the penalty decreases monotonically with parallelism ($0.67$ at 64\,B on
8 CUs, $0.88$ at 512\,B on 40 SMs, $1.00$ at 512\,B on 108 SMs). We report this
as a result: the empirical behaviour converges across four memory architectures without
a predictive model that survives contact with the data, and we prefer to say so
rather than fit a third model post hoc.

\subsection{Fusion works, at small quota}

\begin{figure}[t]
\centering
\includegraphics[width=\columnwidth]{figs/fusion.pdf}
\caption{H200, $Q=64$ queries, real grids. Fusion beats the two-kernel scheme
only when the per-stratum quota is small; labels give winning configurations.}
\label{fig:fusion}
\end{figure}

\begin{table}[t]
\caption{Fused versus two-kernel ASP on H200 (59 configurations, $Q=64$,
$G\in\{32,128,512\}$).}
\label{tab:fusion}
\begin{tabular}{rrrr}
\toprule
$q$ & fused / two-kernel & fused / dense & winning \\
\midrule
\textbf{8}  & \textbf{1.57$\times$} & \textbf{1.71$\times$} (max 4.58) & \textbf{15/15} \\
16 & 0.78$\times$ & 1.27$\times$ & 3/15 \\
32 & 0.68$\times$ & 0.71$\times$ & 0/15 \\
64 & 0.32$\times$ & 0.38$\times$ & 0/14 \\
\bottomrule
\end{tabular}
\end{table}

Table~\ref{tab:fusion} and Figure~\ref{fig:fusion} give the result. Since
$G\cdot q=m$ is fixed, $G$ governs parallelism and $q$ the serial work per team;
the optimum is reached when the per-team quota $q/N_T$ falls to one, i.e.\
$q\approx N_T=8$, giving the design rule $G\approx m/8$. At that point fusion
beats the two-kernel form in 15 of 15 configurations, by $1.57\times$ median.

Reaching this required two diagnoses. The first fused kernel failed with an
illegal memory access, isolated by bisection to partial initialisation of a
shared array whose full extent was later scanned. The second ran but was
$3\times$ slower, with a runtime \emph{exactly linear in $q$} and independent of
summary size---the signature of serialisation, traced to a final merge of
per-team candidate lists executed by a single thread. Removing that merge by
stratifying one level further (each team gathers its own top-$q/N_T$) moved the
result from $0.31\times$ to $0.96\times$, and tuning $G$ to $1.57\times$.

\subsection{End to end: byte savings do not become speedups}

\begin{figure*}[t]
\centering
\includegraphics[width=\textwidth]{figs/e2e.pdf}
\caption{Qwen3-8B on H200. \textbf{Left:} full decode-step latency and the
corresponding byte-traffic ratio. \textbf{Right:} ventilation---dense attention
runs at $84\%$ of peak bandwidth and costs $2\text{--}11$\,ms, while the unfused
ASP selector costs a flat $20\text{--}23$\,ms.}
\label{fig:e2e}
\end{figure*}

\begin{table}[t]
\caption{Qwen3-8B, complete decode step, on two backends.}
\label{tab:e2e}
\small
\setlength{\tabcolsep}{4.5pt}
\begin{tabular}{rrrrrr}
\toprule
context & KV (GB) & dense & ASP & speedup & bytes \\
        &         & (ms)  & (ms) &         & saved \\
\midrule
\multicolumn{6}{l}{\emph{H200, 3.8\,TB/s}}\\
16\,k  &  2.2 & 29.9 & 49.6 & 0.60$\times$ & 13.5$\times$ \\
65\,k  &  9.0 & 29.8 & 49.5 & 0.60$\times$ & 20.9$\times$ \\
131\,k & 18.0 & 30.0 & 49.4 & 0.61$\times$ & 23.0$\times$ \\
262\,k & 36.0 & 42.1 & 54.0 & 0.78$\times$ & 24.2$\times$ \\
524\,k & 72.0 & 72.6 & 89.8 & 0.81$\times$ & 24.9$\times$ \\
\midrule
\multicolumn{6}{l}{\emph{RTX 4090, 0.93\,TB/s}}\\
16\,k  &  2.2 & 113.9 & 212.9 & 0.54$\times$ & 13.5$\times$ \\
33\,k  &  4.5 & 113.6 & 280.3 & 0.41$\times$ & 17.7$\times$ \\
65\,k  &  9.0 & 188.4 & 212.1 & 0.89$\times$ & 20.9$\times$ \\
131\,k & 18.0 & 199.0 & 219.2 & 0.91$\times$ & 23.0$\times$ \\
\bottomrule
\end{tabular}
\end{table}

Table~\ref{tab:e2e} is the paper's central measurement. ASP reads $13$--$25\times$
fewer bytes and is $1.2$--$1.7\times$ slower on the H200 and $1.1$--$2.5\times$
slower on the RTX 4090. That byte count belongs to the harness's selector, whose fine pass re-ranks on the
$r_2=8$ segment summaries of each surviving block and never reads a raw candidate
key, at a block budget of $1.6\%$ of the keys. The two-stage index whose quality is
validated in the long-context validation above instead reads the raw keys of the $2m$
candidates at a $25\%$ budget, and therefore reads $3.9\times$ fewer bytes than
dense ($1.000\times$ the one-stage index at eight bits, $0.877\times$ at four). The
two are different selectors; their byte ratios must not be conflated. At equal bytes the two mechanisms are not equivalent. On SmolLM2-135M at
$T=1024$ and a $25\%$ budget, re-ranking on two segment summaries per block ($520$
units, byte-neutral) costs $+0.0109$\,nat against $+0.0113$ for eight-bit raw keys
at the same $520$ units, but a single block-mean summary ($456$ units, $0.877\times$)
costs $+0.0304$ against $+0.0117$ for four-bit raw keys---$2.6\times$ worse.
Averaging destroys the within-block peak, the same failure mode as the mean-pooled
block key, whereas quantisation preserves the rank. The harness's $r_2=8$ summaries
average eight keys each, i.e.\ the heavy-dilution end of this ladder, so its byte
advantage is not free on a model whose landscape has curvature; on Qwen2.5-0.5B,
where it is flat at this operating point, every variant stays within $0.002$\,nat of
dense. The penalty is regime-dependent rather than intrinsic: it appears only where the
coarse index leaves work to do. At $T=1024$ with the offset at $200{,}000$ keys the
coarse index already costs $+0.013$\,nat and a single block-mean summary is as good
as any other variant ($-0.0051$ against $-0.0050$ for exact raw keys); at a $50\%$
budget it is $-0.0006$ against $+0.0001$. The block-mean summary, the cheapest
mechanism on this ladder, therefore suffices whenever the coarse index is within
about $0.02$\,nat of dense, and loses only when the coarse gap is large ($+0.059$\,nat
at a $25\%$ budget on the same model). The stage-two mechanism should be chosen from
the coarse gap, not from an architectural preference. The budget ladder confirms the prediction. The same comparison at $12.5\%$ and
$6.25\%$ of keys gives dilution penalties of $+0.044$ and $+0.112$\,nat, and at
$6.25\%$ a single block-mean re-rank ($+0.256$) is worse than not re-ranking at all
($+0.244$), while four-bit raw keys ($+0.143$) are the best variant measured. The
harness operates at $1.6\%$ of keys, below this range, so its summary re-rank is
predicted to be penalised at its own operating point; raw-key quantisation is the
mechanism to prefer once the coarse gap exceeds $0.05$\,nat. Tested at the harness's own dilution this is decisive. With $L_b=64$ and $r_2=8$
segments of eight keys---exactly the harness's summaries---the summary re-rank
recovers $10.9\%$ of the stage-two gain, while two-bit raw keys at the same $16$
units recover $76.4\%$ ($+1.280$ against $+0.830$); doubling the number of summaries
to sixteen, four keys each, still recovers only $55.5\%$. More summaries do not
restore the rank that averaging destroys, and quantisation wins at every rung of the
equal-byte ladder ($16$, $8$, $4$ and $2$ units: gaps of $0.09$, $0.33$, $0.36$ and
$0.36$\,nat). The pattern is not specific to this architecture. On Qwen2.5-0.5B at a $1.6\%$
budget the coarse gap is $+0.680$\,nat; two- and four-bit raw keys recover $90\%$ and
$78\%$ of it, two summaries per block recover $40\%$, and a single block-mean summary
is $0.27$\,nat \emph{worse} than not re-ranking at all. At a $6.25\%$ budget on the
same model two summaries per block are the best variant measured ($117\%$) while one
summary is again worse than the coarse index. The governing parameter is the number of
keys averaged per summary: four collapse at sparse budgets on both models, and eight
---the harness's choice---lie beyond that. Precision has a floor on the raw-key route: at the same $1.6\%$ budget, four- and
two-bit keys recover $90.5\%$ and $82.4\%$ of the stage-two gain while a single bit
recovers $-53\%$, i.e. it is worse than not re-ranking at all. Two bits---a $16\times$
compression---are therefore both cheaper and far better than eight-key summaries. The floor does not move with the budget: a single bit is worse than not re-ranking at
$1.6\%$, $6.25\%$ and $25\%$ of keys ($-53\%$, $-52\%$, $-14\%$), while two bits
recover $65$--$82\%$ and four bits $90$--$101\%$ at all three. The same floor appears on the second architecture: on Qwen2.5-0.5B at the same $1.6\%$
budget a single bit costs $+1.71$\,nat against the coarse index's $+0.68$, while two,
four and eight bits recover $53\%$, $78\%$ and $90\%$ of the stage-two gain. Substituting those raw keys for the summaries is therefore strictly dominant: on the
harness's own accounting the byte ratio rises from $13.5$--$23.0\times$ to
$14.8$--$27.3\times$ at two bits---the $m r_2$ term is replaced by $k L_b p$, eight
times smaller---while the recovered stage-two gain rises from $11\%$ to $76\%$. The
summary mechanism is dominated on both axes at once. That dominance is specific to fine re-ranking, and the boundary is the number of keys
per summary. At the harness's block size ($L_b=64$) the mean-summary re-rank is
monotone in dilution and never beats the exact re-rank: $16$, $32$ and $64$ keys per
summary give $+1.333$, $+1.476$ and $+1.829$\,nat against the coarse index's
$+1.355$---the first is neutral, the other two are worse than not re-ranking at
all---while the exact re-rank gives $+0.668$. Shrinking the coarse index's projection
rank also degrades it ($d'=8$, $4$ and $2$ give $+1.355$, $+1.615$ and $+1.603$\,nat),
so the first pass does need its four elements per block. The harness's first pass is a
mean-summary score too, and there it is not merely neutral but catastrophic: at four
elements per block a maximum over four mean summaries of sixteen keys costs
$+3.732$\,nat against the PCA index's $+1.355$, and the gap is unchanged at two and one
element per block ($+3.837$ against $+1.615$; $+3.875$ against $+1.603$). Mean-pooling
attenuates the peak key by the size of the segment, which is fatal when the score must
rank all blocks and only mildly damaging when it must re-order two candidates already
known to contain the peak. The harness uses summaries at both stages, and they buy
nothing at either. Normalising the summaries before scoring does not rescue them: the same
configuration costs $+3.003$\,nat with unit-normalised segment means, still $1.65$\,nat
worse than the index at identical bytes. The gap travels: on Qwen2.5-0.5B the same four-element comparison gives
$+4.187$ against $+1.057$, and unit-normalised summaries still cost $+2.286$ against
$+1.057$. At that block size the fine stage's summaries recover nothing either: with the
candidates chosen by the index, exact re-ranking recovers the whole available gain on
both models ($+0.562$ against $+1.057$ on Qwen2.5-0.5B), while re-ranking by two, four
or one mean summaries of 32, 16 or 64 keys recovers $-16\%$, $-43\%$ and $-75\%$ of it
there --- every granularity worse than not re-ranking at all. Normalising the fine stage's summaries makes it worse rather than better:
the same configurations recover $-66\%$, $-61\%$ and $-50\%$ of the gain on
SmolLM2-135M and $-197\%$, $-135\%$ and $-156\%$ on Qwen2.5-0.5B. The asymmetry with
the coarse stage has a mechanism: there the decision is which block contains the peak,
so the direction suffices and the norm is noise; here every candidate already contains
the peak, so the decision is how much mass it carries, and the norm of a segment mean
is exactly that signal. Normalisation helps a stage that must identify and destroys one
that must weigh. The rule does not depend on the dilution regime: at $L_b=4$, where the summaries
are nearly neutral, normalisation is \emph{more} damaging than at $L_b=64$ --- $-0.91$ and
$-0.94$\,nat on SmolLM2-135M, $-1.38$ and $-1.75$\,nat on Qwen2.5-0.5B. At that block
size the fine stage has almost no headroom anyway (the gain exact re-ranking can recover
falls to $0.007$\,nat on SmolLM2-135M and $0.070$\,nat on Qwen2.5-0.5B), because two
candidates drawn from 254 blocks leave the coarse stage's error in charge. Summaries are not a marginal cost either: they are $32\%$, $41\%$, $49\%$ and $54\%$ of
the harness's selector bytes at $16$k, $32$k, $64$k and $131$k tokens. Replacing both
passes by the architecture validated here---an eight-projection eight-bit index for
ranking and two-bit raw keys of the candidate blocks for re-ranking---cuts the selector
by $30$--$51\%$ and raises the byte ratio from $13.5$--$23.0\times$ to
$19.2$--$47.1\times$, while recovering $76\%$ rather than $11\%$ of the stage-two gain. At the harness's own budget fraction, $1.6\%$ of keys, the contrast is stark. On
SmolLM2-135M at $T=1024$ the coarse index costs $+1.275$\,nat; a four-bit raw-key
re-rank recovers $90.5\%$ of the gap it leaves, and a block-mean re-rank recovers
$15.8\%$ ($+1.241$ against $+1.079$). Across the whole ladder quantisation recovers
$78$--$104\%$ of the stage-two gain at every budget, while the summary mechanisms
collapse from $96\%$ to $35\%$ (two summaries per block) and from $104\%$ to $16\%$
(one summary per block) as the budget falls from $50\%$ to $1.6\%$. We had predicted a $10$--$40\%$ gain; we measure a
$19$--$46\%$ loss. The 4090 reproduces the mechanism but not the shape: its gain
is \emph{not} monotone in context ($0.54$, $0.41$, $0.89$, $0.91$), because the
ASP step is dominated by a flat selector cost while the dense step is not.

The ventilation (Figure~\ref{fig:e2e}, right) identifies the cause completely.
Dense attention costs $2.06$, $4.06$, $6.62$ and $11.34$\,ms at $16$k, $65$k,
$131$k and $262$k---about $84\%$ of the peak-bandwidth floor, i.e.\ already
near-optimal. The unfused ASP selector costs a flat $20$--$23$\,ms: eight
non-fused PyTorch operations per layer, some 290 kernel launches per decode step.
Crossover requires dense attention to exceed the selector overhead. Fitting the
four ventilation points gives $N\approx700$k on the H200---beyond the largest
length we measured, where the speedup is still $0.81\times$.

The same ventilation on the RTX 4090 locates the crossover inside the measured
range. Dense attention there costs $0.161$\,ms per $1{,}000$ tokens of context
(residual $0.8$\,ms) and runs at $80$--$100\%$ of the card's measured peak
bandwidth, while ASP's attention is flat at $24.4$\,ms (residual $1.9$\,ms). The
lines cross at $N\approx146$k tokens, bracketed by measurement: at $131$k ASP is
still slower ($25.3$ against $22.3$\,ms), and at $262$k it is $1.8\times$ faster
($24.2$ against $42.7$\,ms). Because the selector cost is flat, the crossover
scales inversely with dense-attention throughput: the dense slope is $4.3\times$
steeper on the 4090 than on the H200, against a bandwidth ratio of $4.1\times$. A
slower GPU therefore reaches parity \emph{sooner}---the opposite of the intuition
that sparse attention needs fast hardware.

The prediction failed because it assumed a memory-bound decode. This integration
is not: the measured step is $3.2\times$ above its memory floor ($30$\,ms against
$9.4$\,ms at $131$k), the non-attention portion being dominated by framework
overhead. The 4090 makes the point more sharply: its step is $5.6$--$7.3\times$
above the memory floor at $16$k--$131$k, and dense attention is only $2.3\%$,
$6.4\%$ and $11.2\%$ of the step at those lengths---so even a perfect sparse
attention can buy at most a few per cent end-to-end at $16$k. In an engine with CUDA graphs and fused kernels that portion would fall
towards $\approx$5\,ms, attention would dominate, and the byte savings would
matter---provided the selector is fused. That projection is an inference, not a
measurement, and we mark it as such.

\paragraph{An artefact that would have inverted the conclusion.} Our first
ventilation reported a $10.3\times$ ASP speedup on attention. It was wrong: both
paths materialised the expanded KV via \texttt{repeat\_interleave} to emulate
grouped-query attention~\cite{ainslie2023gqa}, inflating the dense baseline by the group factor $g=4$
while barely touching ASP, which gathers a tiny KV. The impossibility was visible
---attention alone ($121$\,ms) exceeded the full step ($30$\,ms). Enabling
GQA-aware SDPA on both paths removed it. We report this because the artefact
favoured our own method, and because the check that caught it (a component
exceeding its container) is cheap and general.

%======================================================================
\subsection{A corrected selector: the maximum, not the mean}
\label{sec:corrected}

An audit of our quality harness, performed after the first version of this paper,
changed our reading of the selector. We report the correction with the same
prominence as the negative result, because it moves the paper's centre of gravity
from \emph{how much} the selector costs to \emph{what} it computes.

\paragraph{The harness was wrong.} Every perplexity measurement in the first
version used sparse masks that omitted the causal bound $j\le p$. The model could
therefore read \emph{future} tokens through the local-window branch. This inverted
the ranking of selectors: the window-only baseline appeared to beat dense
attention, adding blocks appeared to hurt monotonically, and an oracle scoring
blocks by exact attention mass appeared to be the \emph{worst} selector, because
each additional block diluted the leak. The corrected harness carries a
per-position invariance test---the loss on a prefix must equal the loss on the
prefix alone---whose measured discrepancy is $1.6\times10^{-5}$.

\paragraph{A second bug, found by independent reimplementation.}
Reimplementing the sparse path as explicit per-query loops---recomputing
the projections, the block scores and the mask from the weights---revealed
a second defect. The sentinel value $-1$, used when fewer than $m$ blocks
are eligible, was clamped to index $0$ before being scattered into the
selection mask, so block~$0$ was marked selected whenever fewer than $m$
blocks were eligible---that is, for the first $\approx$160 positions of
every chunk. Block~$0$ contains the attention sink, so the defect
\emph{flattered} every sparse selector. Correcting it raises the cost of
sparsity from $0.034$ to $0.215$\,nat and makes the random control
$10\times$ more discriminating ($+1.69$\,nat), while the estimator and
index comparisons below barely move---which is why we report both sets of
numbers. The two implementations agree to $1.2\times10^{-8}$ on the
attention output at the query where they were compared, and a structural
check confirms the eligible-position count ($101$ at $p=100$, $192$ at
$p=300$).

\paragraph{Extending the verification.} We then re-ran the same independent route
on the second model and on the paths the first pass had not touched: Qwen2.5-0.5B
($14$ query heads, $2$ KV heads, hence grouped-query attention), the causal PCA
selector, and rotary embeddings. The reference recomputes the rotations by hand
from the half-split convention, obtains the basis by a singular value
decomposition of the fit keys rather than an eigendecomposition of their
covariance, ranks blocks by an explicit sort rather than \texttt{topk}, and maps
query head $h$ to KV head $h//7$. Over $84$ comparisons ($2$ layers, $3$ query
positions, $14$ heads), the selected block sets and the masks are \emph{identical}
in every case, and the attention outputs agree to $1.4\times10^{-5}$ relative.
This time the audit confirmed the harness rather than breaking it---the outcome a
verification pass must be able to produce, and the reason we kept running it.

\paragraph{Granularity, not the number of blocks, is the dominant knob.} Our
headline configuration ($L_b=32$, $m=2$, window $128$) was one point in a grid we
had not swept. We swept it: $L_b\in\{16,32,64\}$, $m\in\{1,2,4\}$, three corpus
chunks, with the index and the baselines recomputed for every configuration
(Table~\ref{tab:sweep}). Finer blocks are strictly better at equal or lower cost.
At a byte fraction of $31\%$, $L_b=16$ with $m=2$ costs $0.123$\,nat against the
dense model, while $L_b=32$ with $m=1$ costs $0.245$\,nat; and $L_b=16$ with a
single block reads $28\%$ of the keys and still beats $L_b=32$ with two blocks
($0.153$ against $0.214$\,nat). Coarse blocks are far worse: $L_b=64$ costs
$0.36$--$0.39$\,nat even when it reads $75\%$ of the keys. The mechanism is a
boundary effect. With a window of $W$, a block grid of size $L_b$ leaves a
systematic gap of up to $L_b-1$ keys between the window and the nearest eligible
block, so keys just outside the window---often the most relevant distant
ones---are dropped, while the coarse blocks that are admitted bring in irrelevant
keys. The index cost itself shrinks with $L_b$ and $m$ ($+0.058$ down to
$+0.001$\,nat) and is sometimes \emph{negative}: at $m=4$ the PCA index
occasionally beats \emph{maxip}. Granularity is therefore a cheaper knob than
index accuracy, and $L_b$ should be made as small as the scorer can still rank.

\begin{table}[t]
\centering
\caption{Sweep over block size $L_b$ and quota $m$ (SmolLM2-135M, $T=512$,
window $128$, causal PCA index with $d'=8$; means over three corpus chunks).
``keys'' is the fraction $(mL_b+W)/T$ of key positions read; ``index cost'' is
the loss of the PCA index against \emph{maxip} at the same configuration.}
\label{tab:sweep}
\begin{tabular}{rrrrr}
\toprule
$L_b$ & $m$ & keys & vs dense & index cost \\
\midrule
16 & 1 & $28.1\%$ & $-0.153$ & $+0.058$ \\
16 & 2 & $31.2\%$ & $-0.123$ & $+0.026$ \\
16 & 4 & $37.5\%$ & $-0.114$ & $+0.007$ \\
32 & 1 & $31.2\%$ & $-0.245$ & $+0.056$ \\
32 & 2 & $37.5\%$ & $-0.214$ & $+0.011$ \\
32 & 4 & $50.0\%$ & $-0.206$ & $+0.004$ \\
64 & 1 & $37.5\%$ & $-0.385$ & $+0.026$ \\
64 & 2 & $50.0\%$ & $-0.363$ & $+0.007$ \\
64 & 4 & $75.0\%$ & $-0.360$ & $+0.001$ \\
\bottomrule
\end{tabular}
\end{table}

The boundary effect makes a sharp prediction: holding the byte budget fixed and
shrinking the grid should recover quality without spending a byte. It does. With
$m L_b=64$ and $W=128$---so $37.5\%$ of the key positions in every
configuration---the cost against the dense model falls from $0.385$\,nat at
$L_b=64$ to $0.021$\,nat at $L_b=4$, a factor of $18$, while the index cost stays
flat at $0.007$--$0.011$\,nat (Table~\ref{tab:budget}). The growing candidate pool
does not hurt the index: a finer grid also makes the per-block maximum a tighter
proxy for the true maximum, and that gain outweighs having more blocks to rank. At
$L_b=4$ and $m=16$ the sparse model reads $37.5\%$ of the keys and loses
$0.021$\,nat to dense attention---an order of magnitude less than the
configuration we started from.

The effect is not specific to one model. Repeating the constant-budget sweep on
Qwen2.5-0.5B, a different architecture ($14$ query heads, $2$ KV heads), gives
$+0.802$, $+0.346$, $+0.132$, $+0.035$ and $-0.003$\,nat for
$L_b=64,32,16,8,4$: at $L_b=4$ the sparse model matches dense attention while
reading $37.5\%$ of the keys. On this model the index cost is \emph{negative} in
four of the five configurations ($-0.018$ to $-0.005$\,nat): a maximum over
$8$-dimensional projections is not merely a cheaper stand-in for \emph{maxip}, it
is a better scorer. Averaged over chunks the index cost is $\pm0.02$\,nat either
way, so it is best read as free rather than as a fixed penalty.

Pushing the grid further shows where the effect saturates. On Qwen, $L_b=2$ and
$L_b=1$ give $-0.003$ and $-0.004$\,nat, no better than $L_b=4$ ($-0.003$). At
$L_b=1$ the blocks degenerate to single keys, so the selector becomes exact
top-$m$ selection and the measurement is the true per-key oracle at that budget:
the block structure costs nothing once the grid is fine enough. The practical rule
is $L_b\approx4$--$8$, and below it the quota $m$ grows without buying anything.

\paragraph{The window is a red herring; the selection is the mechanism.} Fixing
$L_b=4$ and the total budget, we swept the window against the quota
(Table~\ref{tab:window}). At $192$ keys read, $W=32$ with $m=40$, $W=64$ with
$m=32$ and $W=128$ with $m=16$ all land within $0.008$\,nat of one another
($+0.027$, $+0.029$, $+0.021$): how the budget is split between contiguous and
selected keys barely matters. What matters is that there \emph{is} a selection. A
pure sliding window reading the most recent $256$ keys---half the context, more
keys than any sparse configuration above---loses $1.23$\,nat, two orders of
magnitude worse than reading $192$ keys of which only $16$ are chosen by the
index. Locality alone buys almost nothing; the two-pass selection does the work.

Pushing the budget down sharpens the contrast. At $128$ keys---a quarter of the
context---the two-pass selector costs $0.033$\,nat, while a pure window of exactly
the same size costs $1.68$\,nat: a factor of $50$ at identical bytes. The index's
own cost rises from $0.008$ at $192$ keys to $0.022$ at $128$, because with fewer
blocks admitted each selection error weighs more. The frontier is shallow: going
from $192$ to $128$ keys costs $0.012$\,nat, so a quarter of the context is enough
for the selector to stay within a few hundredths of dense attention.

\paragraph{A third architecture, without rotary embeddings.} GPT-2 has no rotary
embeddings, uses learned absolute positions, and projects queries, keys and values
with a single matrix, so nothing in the mechanism above is specific to the
Llama-style stack. The results carry over. At a fixed $192$ keys the cost against
dense falls from $0.400$\,nat at $L_b=64$ to $0.216$ at $L_b=32$ and $-0.004$ at
$L_b=4$; at $128$ keys the selector costs $0.020$\,nat while a pure window of the
same size costs $1.42$\,nat, a factor of $70$. The PCA index again beats
\emph{maxip} in every configuration ($-0.008$ to $-0.017$\,nat).

\paragraph{Scaling the configuration with the context.} The scale law above held
the configuration fixed and let the fraction of keys read shrink. Scaling the
configuration with the context instead---keeping the fraction at a quarter by
setting $W=T/8$ and $m=T/32$ blocks of $L_b=4$---the cost against dense
\emph{falls}: $+0.041$\,nat at $T=512$, $+0.012$ at $1024$ and $+0.004$ at
$2048$. At the same fraction a pure window costs $+1.76$, $+1.24$ and $+1.08$\,nat,
so the ratio between selection and locality grows from $43\times$ to $299\times$.
The advantage of selection over locality widens with context length, the opposite
of what a utility-threshold argument based on fixed overhead would predict.

\paragraph{Why $L_b\approx4$: a quantitative account.} The boundary gap of a grid
of size $L_b$ holds $(L_b-1)/2$ keys on average, and the attention density just
behind the window turns out to be flat: measured over all layers and heads, the
mass on the $d$-th key behind the window is $7.8\times10^{-4}$ for every $d$ from
$1$ to $60$, with no decay. The trapped mass is therefore $\rho(L_b-1)/2$ with
$\rho\approx7.8\times10^{-4}$, and it costs $\kappa\approx12$--$25$\,nat per unit of
mass: regressing the measured penalty on the measured trapped mass gives
$\kappa=11.7$ with $R^2=0.986$. Repeating the measurement on Qwen2.5-0.5B and
GPT-2 tests whether the two constants transfer. The \emph{shape} does: the measured
penalty differences are proportional to the trapped-mass differences with $R^2$ of
$0.999$ and $0.992$, and the density $\rho$ is nearly invariant across the three
models ($7.5$, $7.8$ and $8.3\times10^{-4}$ per key, a spread under 10\%)---the flat
floor behind the window is a structural property of these attention distributions.
The \emph{scale} does not: $\kappa$ is model-specific, $11.7$ (SmolLM2), $17.9$
(GPT-2) and $33.4$ (Qwen). The rule $\Delta\approx\kappa\rho(L_b-1)/2$ therefore
predicts a model-specific optimum $L_b\lesssim1+2\varepsilon/(\kappa\rho)$, which at
$\varepsilon=0.02$\,nat gives $5.4$, $4.0$ and $2.5$ for the three models---bracketing
the saturation point of $4$ observed on all three, and ordering them correctly. The
same formula explains the scaling result: at a fixed fraction $\rho$ falls as the
context grows, so the optimal $L_b$ stays small and the penalty shrinks.

\paragraph{What sets $\rho$: a sink plus a uniform background.} A natural guess is
that attention outside the window is spread uniformly, giving
$\rho=(1-\mu_w)/(T-W)$. The measured density just behind the window is $2$--$5\times$
below the global average of non-window keys, which first looks like a trough.
Decomposing the out-of-window distribution resolves it: that distribution has an
effective support of only $15$ keys, and a single key holds $0.67$ of it---the
sequence-initial sink, which sits in the first tenth of the positions and carries
$0.70$ of the out-of-window mass. The remaining $0.30$ is a flat background at
$8$--$9\times10^{-4}$ per key. So the background \emph{is} uniform and $\rho$ measures
it; the apparent trough was the background compared against an average inflated by
the sink. This also explains why the granularity penalty is small and why block
selection works at all: the sink lives in the oldest block, which is always eligible
and therefore never trapped in the boundary gap, so the gap only ever takes
background mass---and the background is flat, which is exactly why ranking it is a
near-tie problem. At $W=128$ and $T=512$ the out-of-window mass is $0.58$ of the
total, so this is not a marginal component.

\paragraph{A valid retrieval bench: GPT-2.} The degeneracy is not universal. On GPT-2
the same repeated-span probe engages a genuine retrieval circuit: the loss on the
fourth copy is $2.95$\,nat below the first, the attention on the distant source copy
is $0.437$ against $0.250$ uniform, and a single head at layer 11 puts its entire mass
on it. At a $37.5\%$ read budget the oracle, which selects blocks by true attention
mass, costs $0.007$\,nat; the deployable index costs $0.226$. The obvious
explanation---that the index fails to retain the attention peak---does not survive.
Across six index variants at the same budget, the peak-retention rate has rank
correlation $0.14$ with the loss, while the \emph{retained attention mass} $S$ has
rank correlation $-1.00$: $S=0.857$ costs $0.226$\,nat, $S=0.789$ costs $1.17$,
$S=0.553$ costs $4.64$. That conclusion, however, does not survive a wider sweep.
Replacing the maximum over keys by a log-sum-exp with a tuned temperature---which
interpolates continuously between the maximum and the sum---yields a configuration
that retains \emph{more} mass than the maximum at $25\%$ ($S=0.770$ against
$0.767$) and still loses more ($0.782$ against $0.695$\,nat); and across $18$
configurations spanning three budgets ($25\%$, $37.5\%$, $50\%$) and six
aggregations, the plain maximum is first at every budget. Mass is therefore a good
but not sufficient proxy, because the output is $\sum_k A(k)v_k$: a block can carry
high attention mass and still contribute the wrong value. Both proxies are
imperfect---across the $18$ configurations $S$ and the peak-retention rate have rank
correlations $-0.965$ and $-0.977$ with the loss, and at a fixed $25\%$ budget
$-0.66$ and $-0.83$---so the design rule is not ``maximise a scalar'' but ``prefer
the aggregation that concentrates on the single most aligned key''. On this bench
that is the maximum over projections and over keys, which is also what the
flowing-text regime preferred. Adding value information to the score does not improve on it: reweighting the
block score by the alignment between a block's mean value and the value at the
top-scoring key costs $0.264$\,nat at $37.5\%$ against $0.226$, a random
direction costs $0.433$, and the anti-aligned direction $0.702$---so the value
direction carries real information that a multiplicative factor cannot exploit.
Forcing the selector's own top-scoring block to be kept changes the result not at
all, which locates the damage in the filling of the budget rather than in the
first choice: the attention peak is not the selector's argmax to begin with. Nor does reading less help: thresholding the block score---keeping a block only
if it scores within a factor $\tau$ of the best in its cell---saves blocks only by
removing mass-bearing ones. At $m=24$, raising $\tau$ from $0$ to $0.5$ cuts the
blocks read from $24.0$ to $11.8$ and costs $+0.72$\,nat, whereas truncating the
rank from $24$ to $8$ blocks costs $+0.66$\,nat: at equal bytes the rank dominates
the threshold, because the max-over-keys score is a spike detector and not an
estimator of block mass. The fix is a second stage: keep the coarse index to form a candidate set of $2m$
blocks, then re-rank those candidates with an accurate score. At $37.5\%$ this
brings the penalty from $0.226$ to $0.011$\,nat---the oracle's own $0.007$---and at
$25\%$ from $0.695$ to $0.246$. The selection rule was never the problem; the score
estimator was. The gain is not automatic: re-ranking the candidates with $32$ or
$64$ principal directions helps at the $25\%$ budget ($0.204$ and $0.220$\,nat) but
is six to sixteen times worse than the exact score at $37.5\%$ ($0.069$ and
$0.184$) and catastrophic at $50\%$ ($0.358$, worse than not re-ranking at all),
even though it retains $94.9\%$ of the attention mass and $99.9\%$ of the peak:
mean retained mass is not a sufficient statistic for the loss. Nor is the compact index's problem the basis: replacing the warm-up basis by
the oracle basis---the $32$ dominant singular directions of the current
sequence's own keys---leaves the penalty at $0.166$\,nat at $37.5\%$ and $0.282$
at $50\%$, against $0.011$ and $0.008$ for the exact score. A rank-$32$ score
approximates the value of $q\cdot k$ well (it retains $91.3\%$ of the attention
mass) and still orders the top blocks wrongly: what matters is not how much of
the score the index captures but whether it preserves the ranking of the best
blocks. The compact index is not a sufficient statistic for selection, and the
second stage has to read the keys of its candidates. Whether that second pass can be quantised is model-dependent: on GPT-2, $8$-bit
keys (per-head scale) cost $0.179$\,nat against $0.011$ for the exact score even
though every element errs by only $0.8\%$, whereas on SmolLM2-135M the same
$8$-bit re-ranking is as good as the exact one ($0.042$ against $0.051$\,nat) and on Qwen2.5-0.5B free as well
($0.046$ against $0.048$\,nat)---so
the byte saving cannot be assumed and must be measured per model. The mechanism is a property of the loss landscape, not of the selector:
replacing one of the $16$ selected blocks at random costs $0.167$\,nat on GPT-2
but only $0.010$\,nat on SmolLM2 and none on Qwen2.5-0.5B, and the gap
persists at $k=2$ and $k=4$
($0.286$ against $0.026$, $0.677$ against $0.147$). The GPT-2 bench sits at the
optimum, so any perturbation shows; the SmolLM2 bench sits on a floor of
$0.051$\,nat, where the same perturbation is invisible. The design and the rule both survive at long context: on SmolLM2 at
$T=1024$ with a $25\%$ budget, the coarse index costs $0.059$\,nat, the exact
re-ranking $0.012$ and the $8$-bit re-ranking $0.011$, while re-ranking random
candidates costs $0.694$; the measured $\Delta(1)=0.015$\,nat is flat, so the second stage can be quantised. With eight-bit candidate keys it reads
exactly the same bytes as the one-stage index ($1.000\times$) while being five
times more accurate, and with four-bit keys it reads $0.877\times$; the exact
stage costs $1.246$--$1.330\times$. Quantising is only free where $\Delta(1)$ is
flat: on GPT-2, where it is $0.167$\,nat, eight-bit re-ranking recovers none of
the gain. The plateau is robust: on a second offset and at $T=2048$, $\Delta(1)$ stays
between $0.014$ and $0.018$\,nat, and $4$-bit keys at the second stage cost
within $0.0006$\,nat of the exact re-ranking, so the stage can be built at
four bits. At $T=1024$ and offset $200{,}000$ the two-stage index even beats
dense, by $0.005$\,nat. The result survives a change of architecture: on Qwen2.5-0.5B, at the same
budget, the coarse index costs $0.008$\,nat, the exact re-ranking $0.0007$ (a
factor of $10.7$), and the eight- and four-bit re-rankings $0.002$ and $0.0002$,
all within the $0.002$\,nat noise floor, so the byte-neutral configuration is not
an artefact of one model. A query-free
score, by contrast, fails on both: ranking blocks by the largest key norm, the
family that wins the byte-precision frontier on Qwen3-8B, costs $1.04$, $1.62$ and $1.81$\,nat on GPT-2, SmolLM2 and Qwen2.5-0.5B,
with peak retention falling to $31.6\%$, $23.0\%$ and $30.7\%$. Selection needs a query-dependent score that preserves the order of the
top keys. Budget buys $S$: the loss on the
retrieval task falls through $0.643$, $0.288$, $0.132$ and $-0.001$\,nat at $25\%$,
$37.5\%$, $50\%$ and $75\%$ of keys, and at $75\%$ it matches dense. Projection rank
buys little: at $37.5\%$, $d'=32$ with max scoring reaches $S=0.867$ against $0.856$
for $d'=8$. On retrieval tasks ASP therefore saves about $1.3\times$, not the
$16\times$ it saves on flowing text---the two regimes sit at opposite ends of the
budget axis, but they share one objective, captured attention mass.
The decomposition repeats on both other models. The out-of-window argmax falls in
the first ten positions for $0.97$ (SmolLM2), $0.87$ (GPT-2) and $0.81$ (Qwen) of
queries; the top key holds $0.63$--$0.67$ of the out-of-window mass; the first decile
holds $0.66$--$0.71$; and the remaining deciles are flat at $0.03$--$0.05$. Effective
support varies ($15$, $35$, $21$ keys) and the out-of-window mass itself varies
($0.58$, $0.49$, $0.46$), but the sink-plus-background shape does not.

At longer context the shape gains a third component. Repeating the decomposition at
$T=2048$ the first decile still holds $0.63$ of the out-of-window mass and $0.92$ of
the argmaxes remain in the first ten positions, but the remaining deciles are no
longer flat: their share rises monotonically from $0.019$ at the third decile to
$0.082$ at the last one---a recency gradient peaking immediately behind the window.
The boundary gap therefore sits in the densest part of the background, which is the
worst place to lose keys, while effective support grows from $15$ keys at $T=512$ to
$59$ at $T=2048$ and the sink's share falls from $0.71$ to $0.63$. The measured
penalty nevertheless stays small ($0.004$\,nat at $T=2048$ in the scaled
configuration), so the gradient shifts the mechanism without overturning it.

\paragraph{The constant is not constant in context length.} Repeating the contrast at
a fixed read fraction ($37.5\%$), with the trapped mass measured in the same pass
rather than borrowed from $\rho$, gives at $T=1024$: $\Delta_4=+0.022$,
$\Delta_{16}=+0.036$, $M_4=2.5\times10^{-3}$, $M_{16}=1.2\times10^{-2}$, hence
$\kappa=1.4$; and at $T=2048$: $\Delta_4=-0.055$, $\Delta_{16}=-0.018$,
$M_4=1.6\times10^{-3}$, $M_{16}=8.0\times10^{-3}$, hence $\kappa=5.7$. The estimate
swings by an order of magnitude across context lengths and, more importantly, the
sign of $\Delta$ flips: at $T=2048$ the sparse model \emph{beats} dense by
$0.055$\,nat at the same read fraction. So $\Delta=\kappa M+c$ is a local
linearisation valid within a fixed configuration, not a law with a universal
$\kappa$---a context-dependent benefit of dropping the diffuse background competes
with the trapped-mass cost. The sign flip is itself the useful fact: at long context
the background is mostly noise and removing it helps. The contrast uses oracle
top-$m$ selection, not the deployable index, and one model, so it bounds the
mechanism rather than the system.

\paragraph{The deployable index beats dense at long context.} The contrasts above use
oracle selection. Replacing it with the deployable causal PCA index---base fitted on
the first $256$ tokens, block score
$\max_{k\in B}\max_{i\le 8}\langle q,u_i\rangle\langle k,u_i\rangle$---at the same
fixed read fraction of $37.5\%$ and $L_b=4$ gives, at $T=1024$: dense $13.5125$,
oracle $13.4733$ ($-0.039$), index $13.4768$ ($-0.036$); at $T=2048$: dense
$13.5443$, oracle $13.5179$ ($-0.026$), index $13.4977$ ($-0.047$). Both beat dense,
and at $T=2048$ the index beats the attention-peak oracle. Two caveats bound this.
First, the sign is not robust to how blocks are scored: selecting by summed mass
instead of peak attention gives $+0.022$ at $T=1024$, so aggregation is a first-order
design choice, not a detail. Second, this is loss at a $2.7\times$ KV reduction, not
latency; the selector cost that drives the negative headline is untouched by it.

\paragraph{Where the gain comes from, and what this testbed cannot test.} Sweeping the
read fraction with the deployable index at $T=1024$ (window and quota scaled together,
$L_b=4$) the loss gain over dense grows monotonically as less is read: $-0.036$ at
$37.5\%$, $-0.077$ at $24.7\%$, $-0.090$ at $12.3\%$ and $-0.144$ at $6.0\%$. Read
naively this says sparse is better the sparser it is. Two control probes show why that
reading is unsafe. A synthetic needle of $16$ tokens placed twice gives a dense
attention mass of $0.0395$ on the needle keys, against $20/512=0.039$ for a uniform
distribution; a natural span repeated four times, where induction is demonstrably
active (the loss on the fourth copy is $0.395$\,nat below the first), gives a dense
mass of $0.263$ on the source copy against $0.250$ uniform. In both cases distant
content carries no attention above chance, so there is no concentrated component for a
selector to lose---consistent with the sink-plus-background structure. The local
testbed therefore cannot decide whether sparsity preserves retrieval; that needs a
model or a corpus in which distant attention concentrates.
The same probe on Qwen2.5-0.5B does not repair this: dense induction is absent
there (the loss on the fourth copy is $0.001$\,nat \emph{above} the first) and the
attention on the source copy is $0.204$, \emph{below} the $0.250$ uniform level, while
the nearest copy receives $0.287$---a mild recency tilt, not retrieval. Scaling the
probe model from $135$\,M to $0.5$\,B does not produce a concentrated distant
component.

\begin{table}[t]
\centering
\caption{Window against quota at fixed total budget, $L_b=4$, SmolLM2-135M,
$T=512$, means over three corpus chunks. The last row reads more keys than any
sparse configuration and still loses two orders of magnitude more.}
\label{tab:window}
\begin{tabular}{rrrrr}
\toprule
$W$ & $m$ & keys read & vs dense & index cost \\
\midrule
32 & 40 & 192 & $+0.027$ & $+0.009$ \\
64 & 32 & 192 & $+0.029$ & $+0.011$ \\
128 & 16 & 192 & $+0.021$ & $+0.008$ \\
64 & 48 & 256 & $+0.026$ & $+0.006$ \\
128 & 32 & 256 & $+0.018$ & $+0.005$ \\
256 & 0 & 256 & $+1.234$ & --- \\
\bottomrule
\end{tabular}
\end{table}


\begin{table}[t]
\centering
\caption{Constant byte budget ($mL_b=64$, $W=128$, so $37.5\%$ of key positions
read in every row), SmolLM2-135M, $T=512$, means over three corpus chunks.}
\label{tab:budget}
\begin{tabular}{rrrrr}
\toprule
$L_b$ & $m$ & candidates & vs dense & index cost \\
\midrule
64 & 1 & 8 & $-0.385$ & $+0.026$ \\
32 & 2 & 16 & $-0.215$ & $+0.011$ \\
16 & 4 & 32 & $-0.115$ & $+0.007$ \\
8 & 8 & 64 & $-0.059$ & $+0.008$ \\
4 & 16 & 128 & $-0.021$ & $+0.008$ \\
\bottomrule
\end{tabular}
\end{table}

\paragraph{What the corrected measurements say.} With the mask fixed, the problem
is far better behaved than the first version suggested (Table~\ref{tab:corrected}).
Scoring a block by the maximum inner product over its keys (\emph{maxip}) matches
an exact-mass oracle to within $0.0002$\,nat (three corpus chunks).
$m\in\{1,4,8\}$ selected blocks. The mean-pooled block key that ASP uses ranks
blocks worse than chance on both models. This explains, after the fact, why four
independent refinements of the summary---optimal-weight quadrature, certified
selection, calibrated coresets, adaptive $k$---all failed to improve ranking:
they were refining a statistic that had already destroyed the signal.

\begin{table}[t]
\centering\small
\caption{Corrected selector quality. SmolLM2-135M, $T=512$, window $128$, block
$32$, $m=2$, means over three corpus chunks; only intra-chunk comparisons are
interpretable, since the spread across chunks ($\approx1$\,nat) dwarfs the
differences between selectors. ``Bytes'' is the index per key relative to the full
key vector, at equal per-component precision.}
\label{tab:corrected}
\begin{tabular}{lrrr}
\toprule
selector (score per block) & bytes & loss & $\Delta$ vs \emph{maxip} \\
\midrule
dense attention       & $100\%$   & $2.6108$ & $-0.2145$ \\
\emph{maxip}          & $100\%$   & $2.8253$ & --- \\
exact-mass oracle     & $100\%$   & $2.8255$ & $+0.0002$ \\
PCA $d'=16$, causal   & $25\%$    & $2.8302$ & $+0.0049$ \\
PCA $d'=8$, causal    & $12.5\%$  & $2.8359$ & $+0.0107$ \\
PCA $d'=8$, transductive & $12.5\%$ & $2.8575$ & $+0.0323$ \\
random blocks         & ---       & $4.5155$ & $+1.6903$ \\
\bottomrule
\end{tabular}
\end{table}

\paragraph{A compact index, with no training.} The selector must score every block,
so its own bytes matter. Projecting queries and keys onto $d'=8$ principal
components per head---$12.5\%$ of key bytes---costs $0.0107$\,nat against
\emph{maxip}; at $d'=16$ ($25\%$ of
bytes) the cost falls to $0.0049$ with a causal basis. The same comparison on Qwen2.5-0.5B---a
different architecture, $14$ query heads and $2$ KV heads---gives $+0.0175$\,nat
at $d'=8$ and $-0.0037$\,nat at $d'=16$, the latter within noise of \emph{maxip},
whose own gap to the oracle there is $0.001$\,nat. A random projection at the same
budget fails ($3.52$ versus $3.12$ for chance), so the basis matters. The basis
needs no training. Freezing it on the first $256$ tokens costs $+0.0032$\,nat
against \emph{maxip} on the same chunk---better than recomputing the basis on
the full chunk ($+0.0551$); freezing on $128$ tokens costs $+0.0476$. A basis
estimated on a different domain transfers: evaluated across academic prose,
encyclopedia and source code, cross-domain bases cost $0.007$--$0.075$\,nat
against $1.52$--$2.05$ for random selection. As the candidate pool grows from
$16$ to $32$ to $64$ blocks ($T=512,1024,2048$), the index costs $+0.055$,
$+0.078$ and $+0.132$\,nat, that is $3.0\%$, $4.2\%$ and $7.5\%$ of the gap
between \emph{maxip} and random selection: the index keeps $93$--$97\%$ of the
selectable gain, while the cost of the sparse budget itself stays near
$0.16$\,nat as the context quadruples. Distilling the basis to imitate
\emph{maxip} does \emph{not} help: after correcting for overfitting it remains
$0.019$\,nat behind the training-free PCA basis at the same dimension
($d'=4$), while sitting $0.035$\,nat \emph{ahead} of it in sample: the gap is
overfitting, not capacity.

\paragraph{What survives from the first version.} The latency findings are kernel-
and memory-level measurements and are unaffected by the harness bug: the
irregular gather remains free on datacenter hardware (now across four
architectures), fusion still wins at small quota, and an unfused selector still
costs about $20$--$25$\,ms per decode step on two backends. What
does \emph{not} survive is the inference that the approach is near its quality
ceiling. The ceiling was the estimator, and it is fixable for $12.5\%$ of the key
bytes.

What the corrected harness does show, and the earlier one hid, is that the
sparse budget itself---window $128$ plus two blocks of $32$, $37.5\%$ of
the keys---costs $0.21$\,nat against dense attention. That cost belongs
to the sparsity budget, not to the selector, and it is the quantity a
deployment must trade against the byte savings.

\paragraph{The basis must not see the future.} The harness fit the PCA
basis on all $T$ keys of the chunk, including keys that lie \emph{after} the
query. We audited this by fitting the basis on the first $256$ keys only---a
causal, streaming-compatible choice---and evaluating on the same chunks. The
transductive basis is \emph{worse}, by $0.0216$\,nat on SmolLM2-135M and
$0.0225$ on Qwen2.5-0.5B: letting the index see the future does not help it,
it hurts it, presumably because the chunk-wide key distribution is dominated
by distant tokens. The numbers in the table are therefore the causal ones,
with the transductive row kept as a control; this is also why the frozen
basis of the previous paragraph beat the full-chunk basis.

%======================================================================
\section{Discussion}\label{sec:disc}

\paragraph{Where the effort should go.} The corrected measurements
(Section~\ref{sec:corrected}) reverse one of our earlier priorities. We previously
concluded that making $k$ larger and cheaper was worth more than making the
selector smarter, because the near-tie structure appeared to cap selector quality.
That ceiling was an artefact: with the maximum score per block the selector sits
within $0.0002$\,nat of an exact-mass oracle, and a training-free PCA index reaches
it within $0.032$\,nat at $12.5\%$ of key bytes. The remaining question is
therefore mechanical rather than statistical---whether reading $12.5\%$ of the key
bytes in an index becomes time---which the fused kernel of
Section~\ref{sec:fusion} suggests it can, but which no end-to-end measurement in
this paper establishes.

\paragraph{Sparse attention and speculation conflict.} Under sparse attention the
$\gamma+1$ speculative candidates select different blocks, so verification must
read their union. Measured on real attention, amortisation efficiency falls to
$41.8\%$ at $\gamma=8$ with token-granular selection, and speculative speedup can
vanish entirely ($1.00\times$ at $\alpha=0.7$). Block granularity partially
rescues it ($65.2\%$), which is an argument for block selection that its
proponents do not make.

\paragraph{Should content-defined blocking be in this paper?} We also find that
partitioning the cache by key \emph{geometry} rather than by position doubles the
attention mass captured by the top-8 blocks ($21.6\%\to43.1\%$ without RoPE).
We deliberately exclude it here. It shares the theoretical frame---the radius
$\rho$ of Proposition~\ref{prop:bound} bounds every summary's error---but answers
a different question (\emph{how to partition}, not \emph{how to allocate
precision}), and its validation requires work this paper does not contain: the
gain is measured at a single $k$ and a single block size, and it collapses once
RoPE is applied---from $\times1.99$ to $\times1.23$ ($44.6\%\to54.8\%$), because
rotation partly destroys the geometric clustering the partition exploits.
Clustering pre-RoPE keys while attending post-RoPE recovers only part of it
($\times1.10$), a tension we have characterised but not resolved. Merging the two would dilute both.

%======================================================================
\section{Limitations}

\textbf{Scope of the hardware claims.} Four GPUs, one of them integrated. No
MI300 or Blackwell part was tested; the granularity threshold has no surviving
predictive model, so extrapolation is unwarranted.

\textbf{The end-to-end integration is not a serving engine.} We patch
\texttt{transformers}, not vLLM or SGLang: no continuous batching, no paged KV,
no CUDA graphs. Our central negative result is therefore conditional on an
integration that is itself $3.2\times$ from its memory floor. The fused kernel of
Section~\ref{sec:fusion} was measured \emph{outside} the model; joining the two is
the obvious next step and we have not done it.

\textbf{Context beyond the trained range.} Qwen3-8B has
$\texttt{max\_position\_embeddings}=40960$. Beyond $40$k we measure performance,
not quality; the memory pattern and latency are unchanged but outputs are not
meaningful. Selection quality was measured separately, on real attention, at
in-distribution lengths.

\textbf{Synthetic KV values and simplified summaries.} Decode latency depends on
cache \emph{shape}, not content, so we synthesise KV beyond $131$k; summaries in
the end-to-end path are segment means rather than $k$-center coresets, a weaker
but cheaper construction whose \emph{size}---all that latency depends on---is
identical.

\textbf{Quality of the fused selector.} Fusion's optimum sits at per-stratum
quota one, which costs $2.5$--$3.2$ relative recall points. We measured that cost
but did not measure its effect on any downstream benchmark such as
RULER~\cite{hsieh2024ruler}.

\textbf{Bandit assumptions.} Equation~\eqref{eq:l1} assumes independence between
the two loss sources and Gaussian score noise. The first is a first-order
approximation; the second is contradicted by our own finding that per-query noise
dispersion matters (Section~\ref{sec:law}).

%======================================================================
\section{Conclusion and Future Work}

Two-pass precision allocation reduces selector traffic by $13$--$25\times$ on a
real model, and that reduction does not become a speedup unless two conditions
hold: the selector must be a single fused kernel, and the surrounding engine must
be memory-bound. We measured the first condition being met ($1.57\times$ over the
two-kernel form on H200) and the second being violated ($3.2\times$ from the
memory floor). The irregular gather, which motivated this study, turns out to cost
nothing on datacenter hardware. A corrected quality harness then changed what we
believe the selector is doing: the mean-pooled key that ASP scores with ranks
blocks worse than chance, whereas a maximum over keys projected onto eight
principal components per head recovers the exact-mass oracle to within
$0.032$\,nat at $12.5\%$ of key bytes, with no training.

Four threads remain open. \emph{Fusing the validated kernel into a serving
engine} is the direct continuation and would test the projection of
Section~5.4. \emph{Content-defined blocking} needs its own study, with the
layer-dependence of its gain as the first question. \emph{The second error
mechanism}---the rank structure of estimator error, which a permutation test
shows is actively harmful by $1.2$--$5.2$ points---is characterised but
unexplained. And the \emph{precision-allocation bandit} deserves a proper
optimality theorem: we have a first-order design rule and a measured
rate--distortion curve, not a lower bound.

\bibliographystyle{ACM-Reference-Format}
\bibliography{refs}

\end{document}
